\documentclass[10pt,a4paper]{amsart}
\usepackage[margin=19mm]{geometry}
\usepackage{amsmath,amssymb,mathtools}
\usepackage[scr=rsfs,cal=euler]{mathalfa}
\usepackage{xfrac}
\usepackage[unicode,hidelinks,bookmarks=false]{hyperref}
\newcommand{\ie}{i.e.\ }
\newcommand{\cf}{cf.\ }
\newcommand{\resp}{respectively\ }
\newcommand{\Subsection}{\subsection}
\providecommand{\nobreakdash}{\nobreak\hbox{-}}
\setlength{\emergencystretch}{2em}
\multlinegap0pt
\usepackage{mathtools}
\usepackage{amssymb}
\usepackage{float}
\usepackage{xcolor}
\usepackage{tikz}
\newtheorem{definition}{Definition}
\newtheorem{lemma}{Lemma}
\title{\LARGE \bf Observer-Based Performance-Barrier Event-Triggered Control of $2\times2$ Linear Hyperbolic PDEs}
\author{Eranda Somathilake, Mamadou Diagne}
\date{Source: arXiv:2604.02508v1 (CC BY 4.0)}
\begin{document}
\maketitle

\noindent\emph{Source and licence.} \LARGE \bf Observer-Based Performance-Barrier Event-Triggered Control of $2\times2$ Linear Hyperbolic PDEs. Version arXiv:2604.02508v1 (CC BY 4.0), distributed by arXiv under the Creative Commons Attribution 4.0 International licence, http://creativecommons.org/licenses/by/4.0/, as recorded in the arXiv licence metadata for this exact version and shown on the abstract page https://arxiv.org/abs/2604.02508v1. Full licence notice: \texttt{licenses/arxiv-2604.02508-CC-BY-4.0.txt}.\par\smallskip

\noindent\textit{Editorial scope.} Counted unit: the article's lemma lemma (source0.tex lines 294--296). The article's complete original proof of it is delivered with it (source0.tex lines 297--310, one \texttt{proof} environment of 14 lines). No mathematics is written, repaired or re-derived: the statement and the proof are the article's own lines, in the article's own notation, and no retained line is substituted or re-flowed.  Retained uncounted as the source-local context the proof needs: lines 251--291.  Nothing was omitted from any retained span, so the package carries no omission record.  No retained line contains a citation, so the wrapper carries no bibliography and no bibliography entry.  No retained line contains a figure environment, an image-inclusion command or drawing code, so no image is delivered and the package declares no asset dependency.  Editorial formatting only, in addition: an AMS \texttt{amsart} preamble that restates the article's own declarations and macros used by the retained lines, namely the environments \texttt{definition}, \texttt{lemma} on separate counters, together with the standard packages those lines need.  The retained environments print in this document as Definition 1, Lemma 1 in order of appearance, whereas the article prints the same items with the numbering of its own full document.  The complete native source of the article as distributed in the arXiv e-print for this exact version is bundled unchanged as \texttt{sources/197741/source0.tex}.
\begin{definition}\label{def:etc}
The increasing sequence of times $I$ at which events are generated is defined according to the following rule:
\begin{equation}
    t_{j+1}\coloneqq \inf\{t\geq t_j + \tau\}\vert m(t)<0, j\in\mathbb{N}\}.\label{tj+1}
\end{equation}
    The dynamic variable $m(t)$ is defined as
    \begin{equation}
    \begin{aligned}
        \dot m(t) =& -\eta m(t) + \kappa_1\|\hat \alpha[t]\|^2 + \kappa_2 \|\hat \beta[t]\|^2\\&  + \kappa_3 \hat\alpha(1,t)^2+ \kappa_4 \tilde \beta(0,t)^2 + cW(t),\label{m1}
    \end{aligned}
    \end{equation}
    for $t\in(t_j,t_j+\tau),j\in\mathbb{N}$ and
    \begin{equation}
    \begin{aligned}
        \dot m(t) =& -\eta m(t) -\theta d(t)^2 + \kappa_1 \|\hat \alpha[t]\|^2 + \kappa_2 \|\hat \beta[t]\|^2\\&  + \kappa_3 \hat\alpha(1,t)^2+ \kappa_4 \tilde \beta(0,t)^2 + cW(t),\label{m2}
    \end{aligned}
    \end{equation}
    for $t\in(t_j+\tau,t_{j+1}),j\in\mathbb{N}$. In addition, define $m(0)=0$, $m(t_j)=\omega_0d(t_j^-)^2,j\in\mathbb{N}_{>0}$, and $m(t_j+\tau)=m((t_j+\tau)^-),j\in\mathbb{N}$. The minimum dwell-time $\tau$ is defined such that
    \begin{equation}\label{f}
    \dot f(t) =
        \begin{cases}
            -a_2f(t)^2 - a_1f(t) - a_0&t\in(t_j,t_j+\tau)\\
            0&t\in(t_j+\tau,t_{j+1})
        \end{cases},
    \end{equation}
    with $f(t_j)=\omega_1$, $f((t_j+\tau)^-)=f(t_j+\tau)=\omega_0$\, and 
    \begin{align}
        \tau = \int_{\omega_0}^{\omega_1}\frac{1}{a_2s^2 + a_1s + a_0}ds.\label{tau}
    \end{align}The estimated performance residual, $W(t)$ is defined as
    \begin{equation}
        W(t) = \mathrm{e}^{-\gamma t}\varrho(t) - \hat V(t),\label{w}
    \end{equation}
    where
    \begin{align}
        \varrho(t) =& \max_{s\in[0,t]}\left\{\mathrm{e}^{\gamma s}\hat V(s)\right\},\label{varrho}\\
        V_1(t) =& \int_0^1 \bigg(\frac{A}{\lambda_1(x)}\mathrm{e}^{-\mu\phi_1(x)}\hat \alpha(x,t)^2\notag\\
        &+\frac{B}{\lambda_2(x)}\mathrm{e}^{\mu\phi_2(x)}\hat \beta(x,t)^2\bigg)dx,\label{v1}\\
        \hat V(t) =& V_1(t) + f(t)d(t)^2 + m(t).\label{hat-v}
    \end{align}
    The values $\eta$, $\kappa_1$, $\kappa_2$, $\kappa_3$, $\kappa_4$, $c$, $\omega_1>\omega_0$, and $\gamma\leq \nu$ are positive design parameters. The constants $A$,$B$,$\mu,\theta$, and $\nu$ are to be specified subsequently. 
    \end{definition}

\begin{lemma}
    Consider an increasing sequence of event times $I=\{t_j\}_{j\in\mathbb{N}}$, $t_0=0$, satisfying $\lim_{j\rightarrow\infty}t_j=+\infty$ determined by the triggering rule \eqref{tj+1}-\eqref{hat-v}. Then the estimated performance residual $W(t)$, satisfies $W(t)\geq0 \,\forall t\geq0$. Subsequently, the dynamic variable $m(t)$ governed by \eqref{m1},\eqref{m2} with $m(0)=0$, $m(t_j)\geq 0,j\in\mathbb{N}_{>0}$, $m((t_j+\tau)^-)=m(t_j+\tau),j\in\mathbb{N}$, where $\tau$ is the MDT, it holds that $m(t)\geq0~\forall t\geq0$.
\end{lemma}

\begin{proof}
    Under Definition~\ref{def:etc}, $I$ is an increasing sequence due to the existence of a MDT $\tau$. Therefore, from \eqref{varrho} and \eqref{w}, we have
    \begin{align*}
        W(t) = \mathrm{e}^{-\gamma t}\max_{s\in[0,t]}\left\{\mathrm{e}^{\gamma s}\hat V(s)\right\} - \hat V(t)\geq0,~\forall t\geq0.
    \end{align*}
    Let $m(0)=0$, $m(t_j)\geq 0,j\in\mathbb{N}_{>0}$, $m((t_j+\tau)^-)=m(t_j+\tau),j\in\mathbb{N}$, then from \eqref{m1} we have
    \begin{equation}
    \begin{aligned}
        m(t) = \mathrm{e}^{-\eta(t-t_j)}m(t_j) + \int_{t_j}^t \mathrm{e}^{-\eta(t -\xi)}\big(\kappa_1\|\hat\alpha[\xi]\|^2\\
        \kappa_2 \|\hat \beta[\xi]\|^2  + \kappa_3 \hat\alpha(1,\xi)^2+ \kappa_4 \tilde \beta(0,\xi)^2 + cW(t)\big)d\xi,
    \end{aligned}
    \end{equation}
    $\forall t\in(t_j,t_j+\tau], j\in\mathbb{N}$. Hence $m(t)\geq0~\forall t\in[t_j,t_j+\tau], j\in\mathbb{N}$ Subsequently, for $t\in(t_j+\tau,t_{j+1}), j\in\mathbb{N}$, $m(t)$ evolves as in \eqref{m2} with $t_{j+1}$ selected according to \eqref{tj+1} to ensure $m(t)\geq0,\forall t\in(t_j+\tau,t_{j+1}),j\in\mathbb{N}$. Applying the aforementioned reasoning for each $t_j\in I$ starting with $t_0=0$, it can be shown that $m(t)\geq0, \forall t\geq0$.
\end{proof}

\end{document}
